% 第 14 讲：Tate 上同调与 K^unr
% 源：src/lecture14.txt（230 行，含 2 处页眉）
\section*{第 14 讲\quad Tate 上同调与 $K^{\mathrm{unr}}$}

设 $G$ 是有限群，$X$ 是 $G$-模复形。设 $P\xrightarrow{\ \sim\ }\Z$ 是 $G$-模的投射复形（拟同构）。则
\begin{itemize}[leftmargin=2.2em,itemsep=1pt,topsep=2pt]
\item $X^{hG}:=\operatorname{Hom}_G(P,X)$ 是\textbf{同伦不变量}；
\item $X_{hG}:=P\otimes_{\Z[G]}X$ 是\textbf{同伦余不变商}；
\item $X_{tG}:=\operatorname{hCoker}\bigl(X_{hG}\xrightarrow{\ N\ }X^{hG}\bigr)$ 是 \textbf{Tate 复形}。
\end{itemize}
前两个构造保持拟同构，把无环复形送到无环复形，因为投射复形平坦。此外，回忆记号
\begin{itemize}[leftmargin=2.2em,itemsep=1pt,topsep=2pt]
\item $H^{i}(G,X):=H^{i}(X^{hG})$ 是群上同调；
\item $H_{i}(G,X):=H^{-i}(X_{hG})$ 是群同调；
\item $\widehat{H}^{i}(G,X):=H^{i}(X_{tG})$ 是 Tate 上同调。
\end{itemize}
最后一个构造在 $G$ 为循环群且 $X$ 集中在 $0$ 次时推广了我们先前的构造。

现在让我们考虑模（而非复形）的 Tate 上同调。设 $G$ 作用在 $M$ 上。则对一个阿贝尔群 $N$，给出映射 $N\to M^{G}$ 等同于给出映射 $f\colon N\to M$，使得对一切 $g\in G$ 有 $g\cdot f(x)=f(x)$。对偶地，给出映射 $M_{G}\to N$ 等同于给出映射 $f\colon M\to N$，使得 $f(g\cdot x)=f(x)$（这是因为余不变商是 $M$ 的商，而不变子模是子模）。于是，由于对一切 $g\in G$ 有 $N(g\cdot x)=N(x)$ 与 $g\cdot N(x)=N(x)$，这些泛性质给出交换图表
\[
\begin{CD}
M @>{\mathrm{id}}>> M\\
@VVV @VVV\\
M_{G} @>{N}>> M^{G}
\end{CD}
\]
其中 $N:=\sum_{g}g$，且范映射穿过不变子模与余不变商。注意若 $\#G$ 在 $M$ 中可逆，则范映射 $N$ 是同构。模仿 Tate 上同调的定义，我们得到
\[
M_{G}/N(M_{G})=M_{G}/N(M)=\widehat{H}^{0}(G,M),
\]
可见同调代数实际上比我们「通常的」代数表现更好！

现在问：对 $G$-模复形，范映射 $N$ 是什么？我们有典范复合
\[
X_{hG}=P\otimes_{\Z[G]}X\to\underbrace{\Z\otimes_{\Z[G]}X}_{\text{逐项余不变商}}\xrightarrow{\ N\ }\underbrace{\operatorname{Hom}_{G}(\Z,X)}_{\text{逐项不变子模}}\to\operatorname{Hom}_{G}(P,X)=X^{hG},
\]
其中最后一个映射经由拉回，而范映射逐项地由模上的范映射给出——由先前的构造我们知道它按预期作用（尽管它只在同伦意义下良定义，等等）。注意「逐项不变子模」取的是逐次的「朴素」不变子模，它不保持拟同构；「逐项余不变商」类似。总之，这给出一个映射，我们记之为 $N\colon X_{hG}\to X^{hG}$。

取集中在 $0$ 次的复形（一般地，对有下界的复形），同伦不变量把该复形进一步向右推；类似地，余不变商把它向左推。但 Tate 上同调同时做这两件事，所以结果无界，往往非常混乱。不过在若干较简单的情形可以计算，例如下面这个：

\textbf{命题 14.1.} 设 $M$ 是 $G$-模，看作集中在 $0$ 次的复形。则
\begin{enumerate}[label=(\arabic*),leftmargin=2.6em,itemsep=1pt,topsep=2pt]
\item 若 $i\ge 1$，则 $\widehat{H}^{i}(G,M)=H^{i}(G,M)$；
\item $\widehat{H}^{0}(G,M)=M^{G}/N(M)$；
\item $\widehat{H}^{-1}(G,M)=\operatorname{Ker}(N)/(g-1)M=\operatorname{Ker}\bigl(N\colon M_{G}\to M^{G}\bigr)$；
\item 若 $i\ge 2$，则 $\widehat{H}^{-i}(G,M)=H_{i-1}(G,M)$。\footnote{原文此处作 $H^{i-1}(G,M)$；按证明中的长正合列应为同调群 $H_{i-1}(G,M)$，疑为笔误。}
\end{enumerate}

\textit{证明.} 复合
\[
M_{hG}\xrightarrow{\ N\ }M^{hG}\to M_{tG}=\operatorname{hCoker}(N)
\]
给出上同调长正合列
\[
\cdots\to H^{i}(M_{hG})\to H^{i}(M^{hG})\to\widehat{H}^{i}(G,M)\to H^{i+1}(M_{hG})\to\cdots.
\]
若 $i\ge 1$，则 $H^{i}(M_{hG})$ 与 $H^{i+1}(M_{hG})$ 均消没（因为 $M_{hG}$ 位于非正次数），由正合性得同构 $H^{i}(G,M)\cong\widehat{H}^{i}(G,M)$。至于 $i=0,-1$，此时 $H^{-1}(M^{hG})=0$，给出正合列
\[
0\to\widehat{H}^{-1}(G,M)\to M_{G}\xrightarrow{\ N\ }M^{G}\to\widehat{H}^{0}(G,M)\to 0,
\]
这证明了 (2) 与 (3)。最后，若 $i\ge 2$，则 $H^{-i}(M^{hG})$ 与 $H^{-i+1}(M^{hG})$ 均消没，由正合性得同构
\[
\widehat{H}^{-i}(G,M)\cong H^{-i+1}(M_{hG})=H_{i-1}(G,M),
\]
这证明了 (4)。\hfill$\square$

作为推论，上同调作为 Tate 上同调出现在更高次，但不在零次；类似地，同调也出现，唯独（关键地）在 $0$ 次例外。当然，这一切都依赖于 $P$ 有界这一事实。

\textbf{定理 14.2（上同调类域论主定理）.} 设 $L/K$ 是非阿基米德局部域的扩张，Galois 群 $G$ 有限。则
\[
(L^{\times})_{tG}\simeq(\Z[-2])_{tG}.
\]

虽然这个同构怎么构造并非一目了然，但其实并不太复杂！下一讲我们将把它归结为一个（非常典范的！）消没性陈述。

那么，这个定理究竟意味着什么？取零阶上同调，得到
\[
K^{\times}/N(L^{\times})=\widehat{H}^{0}(G,L^{\times})=H^{0}\bigl((L^{\times})_{tG}\bigr)\simeq H^{0}\bigl((\Z[-2])_{tG}\bigr)=\widehat{H}^{-2}(G,\Z)=H_{1}(G,\Z)\cong G^{\mathrm{ab}},
\]
正如上一讲所证。我们已知这对局部域的 $2$ 次扩张成立，所以这为局部域的希尔伯特符号提供了巨大的推广！

我们现已回忆同构 $H_{1}(G,\Z)\cong G^{\mathrm{ab}}$ 的构造。首先，我们证明了 $i>0$ 时 $H_{i}(G,\Z[G])=0$。事实上，
\[
\Z[G]\otimes_{\Z[G]}P\simeq\Z[G]\otimes_{\Z[G]}\Z=\Z
\]
是拟同构，因为 $P$ 拟同构于 $\Z$。特别地，所有低阶同调群消没，因为 $\Z$ 在 $0$ 次。然后我们构造了如下短正合列：
\[
0\to I_{G}\to\Z[G]\xrightarrow{\ \varepsilon\ }\Z\to 0,
\]
其中 $\varepsilon$ 由 $g\mapsto 1$ 定义，且我们令 $I_{G}:=\ker(\varepsilon)$，它是 $\Z[G]$ 中所有系数之和为零的元素构成的理想。取群同调长正合列，得到
\[
\underbrace{H_{1}(G,\Z[G])}_{=0}\to H_{1}(G,\Z)\to I_{G}/I_{G}^{2}\to\Z\xrightarrow{\ \sim\ }\Z\to 0,
\]
因为一般地 $M_{G}=M/I_{G}M$，故有同构 $H_{1}(G,\Z)\cong I_{G}/I_{G}^{2}$。

最后，我们通过双向构造的映射证明 $I_{G}/I_{G}^{2}\cong G^{\mathrm{ab}}$。我们断言映射
\[
G\longrightarrow I_{G}/I_{G}^{2},\qquad g\mapsto g-1
\]
在模 $I_{G}^{2}$ 意义下是同态（因而诱导出 $G^{\mathrm{ab}}\to I_{G}/I_{G}^{2}$）。为看出这一点，对 $g,h\in G$ 展开
\[
(g-1)(h-1)=gh-g-h+1\in I_{G}^{2},
\]
于是模 $I_{G}^{2}$ 有 $(g-1)+(h-1)\equiv gh-1$；由于 $I_{G}/I_{G}^{2}$ 交换，该映射穿过 $G^{\mathrm{ab}}$。至于逆，我们有复合
\[
I_{G}/I_{G}^{2}\hookrightarrow\Z[G]/I_{G}^{2}\xrightarrow{\ \sum_{g}n_{g}g\ \mapsto\ \prod_{g}g^{n_{g}}\ }G^{\mathrm{ab}},
\]
其中我们把 $G$ 写成乘法群；可以验证这是同态，且是前一映射的逆。

\textbf{例 14.3.} 设 $L/K$ 是局部域的非分歧扩张，故 $G:=\Gal(L/K)=\Z/n\Z$ 循环，且该同构是典范的，Frobenius 元素对应 $1$。回忆 $\widehat{H}^{0}(G,\OK_{L}^{\times})=0$，即 $N\colon\OK_{L}^{\times}\to\OK_{K}^{\times}$ 是满射。我们通过滤过证明了这一点：相继的子商给出范映射 $k_{L}^{\times}\xrightarrow{\ N\ }k_{K}^{\times}$，其余给出迹映射 $k_{L}\xrightarrow{\ T\ }k_{K}$，二者都是满射（例如后者是因为扩张可分）。我们还证明了 Herbrand 商 $\chi(\OK_{L}^{\times})=1$，这推出 $\widehat{H}^{1}(G,\OK_{L}^{\times})=0$，即 $(\OK_{L}^{\times})_{tG}\simeq 0$ 是拟同构。构造短正合列
\[
0\to\OK_{L}^{\times}\to L^{\times}\xrightarrow{\ v\ }\Z\to 0,
\]
其中 $v$ 表示 $L^{\times}$ 上的正规化赋值。取 Tate 上同调便给出拟同构
\[
\Z_{tG}[-2]\simeq(L^{\times})_{tG}\simeq\operatorname{hCoker}\bigl(0\to(L^{\times})_{tG}\bigr)\simeq\operatorname{hCoker}\bigl((\OK_{L}^{\times})_{tG}\to(L^{\times})_{tG}\bigr)\simeq\Z_{tG},
\]
由定理 14.2，且因为 Tate 上同调与锥交换并保持拟同构。于是 $(L^{\times})_{tG}$ 是 $2$-周期的。注意这是典范的，尽管需要选取一个生成元——因为我们可以选 Frobenius 元素（即 $x\mapsto x^{q}$）。

我们现在转向一般局部域扩张的讨论。

\textbf{定义 14.4.} $K^{\mathrm{unr}}$ 是 $K\subseteq\overline{K}$ 中极大非分歧扩张的（$p$ 进）完备化。

\textbf{例 14.5.}
\begin{enumerate}[label=(\arabic*),leftmargin=2.6em,itemsep=2pt,topsep=2pt]
\item 设 $K:=\F_{q}((t))$。$K$ 的第 $n$ 个非分歧扩张是 $K_{n}=\F_{q^{n}}((t))$，故 $K$ 的极大非分歧扩张是
\[
\bigcup_{n\ge 1}\F_{q^{n}}((t))\subseteq\overline{\F}_{q}((t))=K^{\mathrm{unr}}.
\]
\item 若 $K:=\Q_{p}$，则 $K^{\mathrm{unr}}=W(\overline{\F}_{p})[1/p]$，其中 $W(-)$ 表示维特向量环。\footnote{原文此处与下文均作 $\F_{q}$、$\F_{p}$；按文意应为代数闭包 $\overline{\F}_{q}$、$\overline{\F}_{p}$。}
\end{enumerate}

因此 $K^{\mathrm{unr}}$ 的基本结构是一个「局部域」（不是在局部紧的意义上，因为剩余域不有限，而是在它是一个完备离散赋值环的分式域的意义上），其剩余域为 $\overline{\F}_{q}$。

现在，设 $\pi$ 是 $K$ 的素元，它在每个 $K_{n}$（即 $K$ 的 $n$ 次非分歧扩张）中仍是素元，则 $\OK_{K^{\mathrm{unr}}}$ 是 $\bigcup_{n}\OK_{K_{n}}$ 的 $\pi$ 进完备化。于是有短正合列
\[
0\to\OK_{K^{\mathrm{unr}}}^{\times}\to K^{\mathrm{unr},\times}\xrightarrow{\ v\ }\Z\to 0,
\]
其中 $v$ 是 $K^{\mathrm{unr}}$ 上的正规化赋值（故 $v(\pi)=1$）。然后定义
\[
\Gal(K^{\mathrm{unr}}/K):=\operatorname{Aut}_{\mathrm{cts}}(K^{\mathrm{unr}}/K),
\]
后者是 $K^{\mathrm{unr}}$ 的连续 $K$-自同构。记 $k$ 为 $K$ 的剩余域、$\sigma$ 为 $\Gal(\bar k/k)$ 的 Frobenius 元素，则
\[
\Gal(K^{\mathrm{unr}}/K)\cong\Gal(\bar k/k)\cong\widehat{\Z},\qquad\sigma\mapsto 1.
\]

由于
\[
K=\{x\in K^{\mathrm{unr}}:x=\sigma x\},
\]
我们有消解
\[
0\to K\to K^{\mathrm{unr}}\xrightarrow{\ 1-\sigma\ }K^{\mathrm{unr}},
\]
进而有序列
\begin{equation}
0\to K^{\times}\to K^{\mathrm{unr},\times}\xrightarrow{\ x\mapsto x/\sigma x\ }K^{\mathrm{unr},\times}\xrightarrow{\ v\ }\Z\to 0. \tag{14.1}
\end{equation}
注意 $\pi$ 不可能落在中间映射的像中，因为对一切 $x$ 有 $v(x)=v(\sigma x)$。这给出用 $K^{\mathrm{unr},\times}$ 表达 $K^{\times}$ 的公式，它将是我们后续几讲的主要工具。

\textbf{断言 14.6.} 序列 (14.1) 正合。

\textit{证明.} 这等价于说每个 $x\in\OK_{K^{\mathrm{unr}}}^{\times}$ 都能写成 $y/\sigma y$，其中 $y\in\OK_{K^{\mathrm{unr}}}^{\times}$。这归结为证明映射
\[
\OK_{K^{\mathrm{unr}}}^{\times}\to\OK_{K^{\mathrm{unr}}}^{\times},\qquad x\mapsto x/\sigma x
\]
是满射。由极大理想的滤过的完备性（因为 $K^{\mathrm{unr}}$ 按定义完备），只需在伴随分次层面证明这一点。% [原文疑误：原印「第二式是 $\F_q\to\F_q$、$x\mapsto(1-q)x$，并称其『就是恒等映射』」。按伴随分次的实际推导，$\sigma$ 在 $\bar\F_q$ 上就是 Frobenius $x\mapsto x^q$，故该映射应为 $x\mapsto x-x^q$；原式 $(1-q)x$ 只是把它在 char $p$ 下形式地写成「标量 $1-q=1$」。此处径改为 $x-x^q$，满射性由 $\bar\F_q$ 代数闭保证]
这给出映射
\[
\overline{\F}_{q}^{\times}\to\overline{\F}_{q}^{\times},\quad x\mapsto x/x^{q}
\qquad\text{与}\qquad
\overline{\F}_{q}\to\overline{\F}_{q},\quad x\mapsto x-x^{q}.
\]
第一个是满射，因为对任意 $y\in\overline{\F}_{q}^{\times}$ 都可解 $x^{q-1}=1/y$，而 $\overline{\F}_{q}$ 代数闭。第二个也是满射，因为对任意 $a\in\overline{\F}_{q}$ 方程 $T^{q}-T=a$ 在代数闭域中总有解。\hfill$\square$
